Este capítulo describe JobAI: su arquitectura, el recorrido de una consulta en la configuración de referencia, que es la que se evalúa, y las dos iteraciones de diseño de este trabajo. Los detalles de implementación están en el Anexo~\ref{anx:artefacto}, y las fallas que orientaron las iteraciones, en la Sección~\ref{sec:res-fallas}.

\section{Arquitectura}
\label{sec:art-arquitectura}

JobAI es una plataforma de agregación de ofertas de empleo remoto cuyo componente central es un copiloto conversacional basado en RAG. Se organiza en tres capas (Figura~\ref{fig:art-arquitectura}): un cliente web en Next.js, que atiende la autenticación, la navegación del catálogo y el panel conversacional; un servidor en NestJS, que orquesta la ingesta, la recuperación y la generación, y una capa de datos con PostgreSQL como almacén relacional y un almacén vectorial que, desde la Iteración~1, puede ser pgvector o Qdrant. Un servidor propio, ejecutado localmente sobre una GPU, sirve el modelo de \textit{embeddings} y el de reordenamiento; los modelos de lenguaje se consumen a través de un proveedor externo, Groq. Los módulos del servidor se describen en el Anexo~\ref{anx:artefacto}.

El capítulo describe la \emph{configuración de referencia} de JobAI, definida en el repositorio junto con el código. La Tabla~\ref{tab:art-config} la compara con las otras dos que distingue la Sección~\ref{sec:met-sistema}: la \emph{desplegada} en el servidor público, que usa el mismo código con otro modelo de \textit{embeddings} y sin reordenamiento, y la \emph{evaluada}, que coincide con la de referencia salvo en el modelo con que se analizó el corpus congelado, y registra los modelos y las fechas de cada corrida. Las tres comparten la primera etapa híbrida y los filtros de la Iteración~2.

\begin{figure}[htbp!]
    \centering
    \caption{Arquitectura de JobAI}
    \label{fig:art-arquitectura}
    \begin{tikzpicture}[
        box/.style={draw, rounded corners, align=center, minimum height=1cm, text width=3.2cm, font=\small},
        store/.style={draw, cylinder, shape border rotate=90, aspect=0.15, align=center, minimum height=1.2cm, text width=2.4cm, font=\small},
        arr/.style={->, >=stealth}
    ]
        \node[box, text width=2.2cm] (fe) at (0,0) {Cliente web\\(Next.js)};
        \node[box, text width=2.2cm] (be) at (5.4,0) {Servidor\\(NestJS)};
        \node[box] (models) at (10.2,2.1) {Modelos de lenguaje\\(Groq)};
        \node[box] (local) at (10.2,0) {\textit{Embeddings} y reordenamiento (servidor local, GPU)};
        \node[box] (sources) at (10.2,-2.1) {Fuentes de ofertas\\(Remotive, Jobicy, Findwork)};
        \node[store] (pg) at (3.8,-3.2) {PostgreSQL\\(+ pgvector)};
        \node[store] (qd) at (7.0,-3.2) {Qdrant};
        \draw[arr] (fe) -- node[above, font=\scriptsize] {REST / NDJSON} (be);
        \draw[arr] (be) -- (models);
        \draw[arr] (be) -- (local);
        \draw[arr] (sources) -- (be);
        \draw[arr] (be) -- (pg);
        \draw[arr, dashed] (be) -- (qd);
    \end{tikzpicture}
\end{figure}

\section{Pipeline de consulta}
\label{sec:art-consulta}

La Figura~\ref{fig:art-pipeline} resume el recorrido de una consulta y la Tabla~\ref{tab:art-config}, los modelos y parámetros de las tres configuraciones.

\begin{figure}[htbp!]
    \centering
    \caption{Pipeline de consulta de JobAI}
    \label{fig:art-pipeline}
    \begin{tikzpicture}[
        stage/.style={draw, rounded corners, align=center, text width=4.2cm, minimum height=0.9cm, font=\small},
        half/.style={stage, text width=2.9cm},
        decision/.style={draw, diamond, aspect=2.2, align=center, inner sep=1pt, font=\small},
        arr/.style={->, >=stealth}
    ]
        \node[stage] (q) at (0,0) {Consulta del usuario};
        \node[decision] (cls) at (0,-1.6) {Clasificación};
        \node[stage] (rag) at (-3.6,-3.2) {Ruta de recuperación\\(con restricciones como filtro)};
        \node[stage] (agg) at (3.6,-3.2) {Ruta de agregación\\(plantilla SQL)};
        \node[half] (dense) at (-5.5,-4.8) {Búsqueda densa: 50 fragmentos $\rightarrow$ ofertas};
        \node[half] (bm25) at (-1.7,-4.8) {BM25:\\50 ofertas};
        \node[stage] (rrf) at (-3.6,-6.4) {Fusión RRF\\(20 candidatas)};
        \node[stage] (rr) at (-3.6,-7.8) {Reordenamiento\\(5 ofertas)};
        \node[stage] (ctx) at (-3.6,-9.2) {Contexto por oferta + campos estructurados};
        \node[stage] (sum) at (3.6,-4.8) {Resumen de la agregación};
        \node[stage] (llm) at (0,-10.8) {Prompt + generación (LLM)};
        \node[stage] (out) at (0,-12.2) {Respuesta incremental con fuentes};
        \draw[arr] (q) -- (cls);
        \draw[arr] (cls) -| node[above, pos=0.25, font=\scriptsize] {recuperación} (rag);
        \draw[arr] (cls) -| node[above, pos=0.25, font=\scriptsize] {agregación} (agg);
        \draw[arr] (rag) -- (dense);
        \draw[arr] (rag) -- (bm25);
        \draw[arr] (dense) -- (rrf);
        \draw[arr] (bm25) -- (rrf);
        \draw[arr] (rrf) -- (rr);
        \draw[arr] (rr) -- (ctx);
        \draw[arr] (agg) -- (sum);
        \draw[arr] (ctx) |- (llm);
        \draw[arr] (sum) |- (llm);
        \draw[arr] (llm) -- (out);
    \end{tikzpicture}
\end{figure}

\begin{table}[htbp!]
    \centering
    \caption{Configuraciones desplegada, de referencia y evaluada de JobAI}
    \label{tab:art-config}
    \footnotesize
    \begin{tabular}{|>{\raggedright\arraybackslash}p{2.6cm}|>{\raggedright\arraybackslash}p{3.3cm}|>{\raggedright\arraybackslash}p{3.6cm}|>{\raggedright\arraybackslash}p{3.8cm}|}
        \hline
        \textbf{Componente} & \textbf{Desplegada} & \textbf{De referencia} & \textbf{Evaluada} \\
        \hline
        Modelo de \textit{embeddings} & \texttt{gemini-\allowbreak embedding-\allowbreak 001}, proveedor externo, 768 dimensiones & multilingual-e5-base \citep{Wang2024mE5}, 768 dimensiones, servidor local; prefijos \texttt{query:} y \texttt{passage:} & Igual que la de referencia, en la GPU del equipo de evaluación; e5-base-v2 \citep{Wang2022E5} se conserva para medir el cambio de modelo \\
        \hline
        Almacén vectorial & pgvector & pgvector (valor por defecto) & Qdrant; pgvector para comparar (PE1 y PE4) \\
        \hline
        Índice vectorial & \multicolumn{3}{>{\raggedright\arraybackslash}p{11.5cm}|}{HNSW con $m=16$ y $ef_{construction}=128$ en pgvector y en Qdrant; $ef_{search}=40$} \\
        \hline
        Restricciones & \multicolumn{3}{>{\raggedright\arraybackslash}p{11.5cm}|}{Extraídas por el clasificador (salario, modalidad remota, tipo de contrato) y aplicadas como filtro dentro de la búsqueda densa y de BM25} \\
        \hline
        Primera etapa & \multicolumn{3}{>{\raggedright\arraybackslash}p{11.5cm}|}{50 fragmentos más similares por coseno, agrupados por oferta, y 50 ofertas por BM25, fusionadas con RRF ($k = 60$)} \\
        \hline
        Reordenamiento & Desactivado; se conservan las 5 primeras de la fusión & Qwen3-Reranker-0.6B \citep{Zhang2025Qwen3} sobre las 20 primeras candidatas; se conservan 5 ofertas & Igual que la de referencia \\
        \hline
        Compuerta de sin resultados & No aplica: depende del reordenamiento & Umbral de 0,134 sobre el mejor puntaje & Igual; las métricas de recuperación se calculan sin ella \\
        \hline
        Análisis estructurado (ingesta) & \texttt{gpt-oss-120b} (Groq) & \texttt{gpt-oss-120b} (Groq) & Corpus analizado con Llama 3.1 8B, local (Anexo~\ref{anx:corpus}) \\
        \hline
        Fragmentación & \multicolumn{3}{>{\raggedright\arraybackslash}p{11.5cm}|}{Estructurada, en cuatro tipos: identidad; requisitos; responsabilidades y habilidades; beneficios y salario. Ventana deslizante de 800 caracteres con solapamiento de 200 cuando el análisis estructurado falla} \\
        \hline
        Clasificación y generación & \multicolumn{3}{>{\raggedright\arraybackslash}p{11.5cm}|}{Patrones \textit{regex} y \texttt{qwen/qwen3.8-27b} (Groq) restringido a un catálogo cerrado de plantillas SQL; el mismo modelo genera la respuesta, con temperatura 0,2 (0 durante la evaluación)} \\
        \hline
        \multicolumn{4}{|l|}{\textit{Solo en la evaluación}} \\
        \hline
        Jueces & \multicolumn{3}{>{\raggedright\arraybackslash}p{11.5cm}|}{GPT-6 Sol sin razonamiento y con temperatura 0, para la relevancia y para RAGAS; Llama 3.1 8B, local, solo durante el desarrollo del protocolo} \\
        \hline
        Fechas de las corridas (2026) & \multicolumn{3}{>{\raggedright\arraybackslash}p{11.5cm}|}{Recuperación: 14 de octubre y 16--17 de noviembre (sistemas complementarios); generación: 14 de octubre; latencia: 14 y 15 de octubre y 1 de noviembre; juicios de relevancia: 13--14 y 16 de noviembre; RAGAS: 14 y 16 de noviembre} \\
        \hline
        Equipo & \multicolumn{3}{>{\raggedright\arraybackslash}p{11.5cm}|}{AMD Ryzen 5 3600 (6 núcleos, 12 hilos), 16 GB de RAM y GPU NVIDIA GeForce RTX 3060 de 12 GB, que ejecuta el modelo de \textit{embeddings} y el de reordenamiento} \\
        \hline
    \end{tabular}
\end{table}

Solo tres decisiones de la configuración de referencia se tomaron con datos: el modelo de reordenamiento, el número de candidatas que recibe y el umbral de la compuerta, elegidos con el conjunto de desarrollo y una regla fijada antes de medir (Sección~\ref{sec:art-rerank}). Los demás parámetros no se ajustaron y su valor tiene otro origen: los 50 fragmentos de la búsqueda densa y las 50 ofertas de BM25 se fijaron al diseñar la Iteración~2; $k = 60$ es el valor habitual de RRF \citep{Cormack2009}; las cinco ofertas del contexto vienen del diseño anterior a este trabajo; $ef_{search} = 40$ es el valor por defecto de pgvector, que se usa también en Qdrant, y cada oferta se recorta a 2.000 caracteres antes del reordenamiento para no exceder el límite de 512 \textit{tokens} del modelo. Ninguno de ellos se evaluó, y otros valores podrían dar resultados distintos.

\label{sec:art-ingesta}La ingesta es un proceso programado que consulta las tres fuentes públicas cada seis horas y puede desactivarse, como se hizo para congelar el corpus de evaluación. Esa copia congelada, con las 1.068 ofertas, sus fragmentos y sus vectores, está publicada junto con su manifiesto en \url{https://github.com/esteban0406/RAG-JobPosting/releases/tag/corpus-2026-10-14}. Deduplica las ofertas, extrae de cada una con un modelo de lenguaje sus campos estructurados (resumen, salario, responsabilidades, requisitos, beneficios, habilidades, tipo de contrato y modalidad), normaliza los salarios a dólares anuales y la divide en fragmentos \emph{tipados}: identidad, requisitos, responsabilidades y habilidades, y beneficios y salario (Figura~\ref{fig:art-ingesta}). La fragmentación tipada es anterior a este trabajo y queda fuera de su alcance (Sección~\ref{sec:alcance}): evaluarla exigiría volver a fragmentar y codificar el corpus y rehacer los juicios de relevancia, y se deja como trabajo futuro. La configuración de referencia usa \texttt{gpt-oss-120b} para el análisis estructurado, pero el corpus de evaluación se procesó con Llama 3.1 8B ejecutado localmente (Anexo~\ref{anx:corpus}).

\begin{figure}[htbp!]
    \centering
    \caption{Flujo de la ingesta de ofertas}
    \label{fig:art-ingesta}
    \begin{tikzpicture}[
        stage/.style={draw, rounded corners, align=center, text width=2.9cm, minimum height=1.25cm, font=\scriptsize},
        store/.style={draw, cylinder, shape border rotate=90, aspect=0.15, align=center, text width=2.6cm, minimum height=1.0cm, font=\scriptsize},
        arr/.style={->, >=stealth}
    ]
        \node[stage] (src) at (0,0) {\textbf{Fuentes públicas}\\Remotive, Jobicy y Findwork, cada seis horas};
        \node[stage] (dup) at (3.6,0) {\textbf{Deduplicación}\\resumen SHA-256 de título, empresa y descripción};
        \node[stage] (llm) at (7.2,0) {\textbf{Análisis estructurado (LLM)}\\resumen, salario, requisitos, beneficios, habilidades, contrato, modalidad};
        \node[stage] (sal) at (10.8,0) {\textbf{Normalización de salarios}\\dólares anuales; fuera de rango, no informado};
        \node[stage] (chk) at (10.8,-2.9) {\textbf{Fragmentación tipada}\\identidad; requisitos;\\responsabilidades y habilidades;\\beneficios y salario};
        \node[stage] (emb) at (7.2,-2.9) {\textbf{\textit{Embeddings}}\\multilingual-e5-base, un vector por fragmento};
        \node[store] (pg) at (2.4,-1.9) {PostgreSQL: texto, atributos y pgvector};
        \node[store] (qd) at (2.4,-4.0) {Qdrant: vectores y atributos};
        \node[font=\scriptsize, align=center, text width=3.2cm] at (10.8,-4.75) {si el análisis falla: ventana deslizante de 800 caracteres};
        \draw[arr] (src) -- (dup);
        \draw[arr] (dup) -- (llm);
        \draw[arr] (llm) -- (sal);
        \draw[arr] (sal) -- (chk);
        \draw[arr] (chk) -- (emb);
        \draw[arr] (emb) -- (pg);
        \draw[arr, dashed] (emb) -- (qd);
    \end{tikzpicture}
\end{figure}


\label{sec:art-clasificacion}En la clasificación de la consulta, patrones \textit{regex} detectan señales de recuperación, de agregación y de restricción, en inglés y en español. Una consulta con señales de recuperación y sin las otras va directamente a la recuperación, sin invocar ningún modelo de lenguaje. En los demás casos, un clasificador basado en LLM, a temperatura 0, elige entre recuperación, agregación o una ruta híbrida y, según el caso, selecciona una de 24 plantillas SQL parametrizadas o extrae las restricciones de la consulta: salario mínimo o máximo, modalidad remota y tipo de contrato. El modelo nunca genera SQL libre. Tres reglas deterministas corrigen la salida del clasificador cuando su etiqueta y su plantilla se contradicen (Anexo~\ref{anx:clasificacion}). El título del cargo y las ubicaciones distintas de \textit{remoto} no se convierten en filtros: exigir su coincidencia literal reproduciría la limitación léxica que el pipeline busca superar.

\label{sec:art-recuperacion}En la recuperación, la búsqueda densa recupera los 50 fragmentos más similares que cumplen el filtro y los agrupa por oferta, con la similitud del mejor fragmento; si el usuario cargó un currículo, antepone a la consulta un perfil breve. En paralelo, BM25 \citep{Robertson2009} busca las 50 ofertas que mejor coinciden con el texto de la consulta bajo el mismo filtro. Las dos listas se combinan mediante fusión por rangos recíprocos (RRF) \citep{Cormack2009}, que suma $1/(k + r)$ por cada lista en la que la oferta aparece en la posición $r$, con $k = 60$, sin necesidad de calibrar entre sí los puntajes. Las 20 primeras ofertas pasan por un modelo de reordenamiento de tipo \textit{cross-encoder}, Qwen3-Reranker-0.6B \citep{Zhang2025Qwen3}, que lee juntas la consulta y la oferta \citep{Nogueira2019} y estima la probabilidad de que la oferta sea relevante; se conservan las cinco mejores. El filtro se aplica dentro de cada motor, durante la búsqueda; si ninguna oferta lo cumple, la búsqueda se repite sin él y el prompt pide advertirlo al usuario. Una compuerta de \emph{sin resultados} responde que no hay ofertas relevantes si la mejor oferta reordenada tiene una probabilidad menor que 0,134, un umbral calibrado con consultas de desarrollo (Sección~\ref{sec:art-iter-recuperacion}). Si el servidor de reordenamiento no responde, se conserva el orden de la fusión.

Para generar la respuesta, por cada una de las cinco ofertas se arma un bloque de contexto con su encabezado y el texto de sus fragmentos. Si la consulta pregunta por un campo, como el salario o los requisitos, que no aparece en los fragmentos recuperados, se añade desde los datos estructurados de la oferta, para que el modelo no lo complete con información plausible pero inexistente. El prompt instruye al modelo a responder \emph{únicamente} a partir de las ofertas, a indicar cuándo no contienen la información solicitada, a no inventar detalles y a responder en el idioma de la pregunta. El historial reciente de la conversación, hasta 12 mensajes, se incluye solo en la generación; la recuperación usa la consulta actual. En la ruta híbrida, la recuperación y la agregación se ejecutan en paralelo y se combinan en un único prompt.

\section{Iteración 1: corrección}
\label{sec:art-vectorial}

\label{sec:art-antecedentes}El historial de versiones de JobAI documenta tres iteraciones de diseño y evaluación anteriores a este trabajo, hechas en abril de 2026 sobre un corpus distinto, de unas 3.200 ofertas, y evaluadas con 15 consultas y un modelo de lenguaje como juez (Tabla~\ref{tab:met-iteraciones}, en el Anexo~\ref{anx:protocolo}). En ellas se adoptó la fragmentación estructurada, que este trabajo no vuelve a evaluar (Sección~\ref{sec:alcance}). Su valor como evidencia es limitado: el conjunto de consultas es pequeño y se diseñó a la vista del corpus, no hay un sistema de referencia, el juez nunca se validó contra juicios humanos y el corpus ya no es el que la plataforma sirve; la misma configuración obtuvo una precisión de 90,7\,\% o de 98,7\,\% según el modelo que actuó como juez. El protocolo del Capítulo~\ref{cap:metodologia} responde a esas cuatro limitaciones.

La primera iteración partió de una revisión del sistema frente a su documentación, hecha al preparar la evaluación, para que lo medido correspondiera a lo que el sistema dice hacer. Corrigió fallas que reúne la Sección~\ref{sec:res-fallas}, la mayoría en la ingesta, la agregación y la búsqueda vectorial, y rediseñó el almacenamiento vectorial para que pgvector y Qdrant pudieran compararse en igualdad de condiciones.

En su versión inicial, JobAI almacenaba los vectores en una columna de PostgreSQL mediante la extensión pgvector \citep{Pgvector}, con un índice IVFFlat creado en la migración inicial, sobre una tabla todavía vacía. La consulta de similitud ordenaba los resultados por la expresión derivada $1 - d(\mathbf{q}, \mathbf{v})$, y el planificador de PostgreSQL solo puede servir una búsqueda de vecinos cercanos desde el índice cuando la consulta ordena directamente por el operador de distancia. La inspección con \texttt{EXPLAIN ANALYZE} sobre el corpus de evaluación (Tabla~\ref{tab:art-explain}) confirmó que la consulta original nunca usaba el índice. También mostró que, con el índice HNSW y la consulta corregida, el planificador sigue prefiriendo el barrido secuencial a esta escala, aunque el recorrido del índice es 12 veces más rápido; por eso la búsqueda admite un modo que fuerza el índice.

\begin{table}[htbp!]
    \centering
    \caption{Plan de ejecución de la búsqueda de 15 vecinos en pgvector sobre el corpus de evaluación}
    \label{tab:art-explain}
    \small
    \begin{tabular}{|p{5.2cm}|p{2.2cm}|p{4.6cm}|r|}
        \hline
        \textbf{Consulta} & \textbf{Índice} & \textbf{Plan elegido} & \textbf{Tiempo} \\
        \hline
        Original (orden por $1-d$) & IVFFlat & Barrido secuencial + ordenamiento & 25,7 ms \\
        \hline
        Reescrita (orden por $d$) & IVFFlat & Recorrido del índice & 2,8 ms \\
        \hline
        Reescrita (orden por $d$) & HNSW & Barrido secuencial + ordenamiento & 24,8 ms \\
        \hline
        Reescrita, índice forzado & HNSW & Recorrido del índice & 2,0 ms \\
        \hline
    \end{tabular}
\end{table}

Este diagnóstico cambia la motivación de la migración. La lentitud observada en pgvector no se debía al motor, sino a una consulta que no usaba el índice. Por eso la migración a Qdrant no se planteó como una mejora supuesta de antemano, sino como una alternativa que se evalúa (PE4).

La migración hizo que los módulos de ingesta, RAG y gestión de ofertas pasaran a depender de una interfaz única de almacén vectorial, con dos implementaciones seleccionables por configuración. En pgvector, la consulta se reescribió para ordenar por el operador de distancia, el índice IVFFlat se reemplazó por HNSW \citep{Malkov2020} y las búsquedas con filtros activan los barridos iterativos del índice. En Qdrant \citep{QdrantDocs}, la colección usa los mismos parámetros HNSW y guarda como carga útil los atributos que filtran la búsqueda. PostgreSQL sigue siendo la fuente de verdad del texto y de la relación con las ofertas, y un comando idempotente copia los vectores a Qdrant y verifica los conteos (Figura~\ref{fig:art-iter1}; detalles en el Anexo~\ref{anx:migracion}). Sobre el corpus de evaluación, la búsqueda exacta en ambos motores coincidió en el 99,85\,\% de los 10 primeros resultados de 200 consultas con los vectores de e5-base-v2, y en el 99,75\,\% tras recodificar el corpus con el modelo multilingüe; las diferencias son empates entre fragmentos de texto idéntico. Los dos almacenes contienen, por tanto, los mismos datos y usan la misma métrica.

\begin{figure}[htbp!]
    \centering
    \caption{Almacenamiento vectorial antes y después de la Iteración~1}
    \label{fig:art-iter1}
    \begin{tikzpicture}[
        box/.style={draw, rounded corners, align=center, text width=2.6cm, minimum height=0.9cm, font=\scriptsize, inner sep=4pt},
        wide/.style={box, text width=5.2cm},
        store/.style={draw, cylinder, shape border rotate=90, aspect=0.15, align=center, text width=2.6cm, minimum height=1.0cm, font=\scriptsize},
        note/.style={align=left, text width=5.6cm, font=\scriptsize},
        new/.style={fill=AzulInstitucional!10},
        arr/.style={->, >=stealth}
    ]
        \node[font=\small\bfseries] at (3.0,3.5) {Antes};
        \node[wide] (s0) at (3.0,2.6) {Módulos de ingesta, RAG y gestión de ofertas};
        \node[store] (p0) at (3.0,0.9) {PostgreSQL + pgvector};
        \draw[arr] (s0) -- node[right, font=\scriptsize, xshift=2pt] {acceso directo} (p0);
        \node[note] at (3.0,-0.9) {Índice IVFFlat creado sobre la tabla todavía vacía.\\Consulta ordenada por $1 - d(\mathbf{q},\mathbf{v})$: el planificador hace un barrido secuencial y nunca usa el índice.};
        \draw[gray!60] (6.4,-2.9) -- (6.4,3.7);
        \node[font=\small\bfseries] at (10.0,3.5) {Después};
        \node[wide] (s1) at (10.0,2.6) {Módulos de ingesta, RAG y gestión de ofertas};
        \node[wide, new] (if) at (10.0,1.1) {Interfaz única de almacén vectorial\\(implementación elegida por configuración)};
        \node[box, new] (pg) at (8.1,-0.75) {\textbf{pgvector}\\HNSW; orden por $d(\mathbf{q},\mathbf{v})$; barridos iterativos con filtros};
        \node[box, new] (qd) at (11.9,-0.75) {\textbf{Qdrant}\\HNSW con los mismos parámetros; atributos como carga útil};
        \draw[arr] (s1) -- (if);
        \draw[arr] (if) -- (pg);
        \draw[arr] (if) -- (qd);
        \draw[arr, dashed] (pg) -- node[above, font=\scriptsize] {copia} (qd);
        \node[note, text width=6.6cm] at (10.0,-2.4) {PostgreSQL sigue siendo la fuente de verdad del texto; un comando idempotente copia los vectores a Qdrant y verifica los conteos.};
    \end{tikzpicture}
\end{figure}


\section{Iteración 2: mejora}
\label{sec:art-iter-recuperacion}

Una primera corrida completa de la evaluación, con la configuración de la Iteración~1, y una revisión del pipeline frente a la literatura mostraron cuatro brechas. La primera era que las restricciones de la consulta no llegaban a la búsqueda vectorial (falla A4): cuando el clasificador sí las extraía, la plataforma las resolvía con un filtro SQL por subcadena del título, la misma búsqueda léxica que el trabajo pretende superar, y el \textit{benchmark} de búsqueda filtrada (PE4) medía una capacidad que la aplicación no usaba. La segunda, que el pipeline era un RAG \emph{ingenuo} en la clasificación de \citet{Gao2023}, con una única búsqueda densa y sin reordenamiento ni recuperación híbrida. Las otras dos tenían que ver con el idioma: las consultas en español recibían respuestas en inglés (falla A6), y el modelo de \textit{embeddings}, e5-base-v2 \citep{Wang2022E5}, estaba entrenado solo en inglés, de modo que nada garantizaba que acercara una consulta en español a la oferta equivalente en inglés.

La Figura~\ref{fig:art-iter2} compara el pipeline antes y después de la iteración, y la Tabla~\ref{tab:art-cambios} resume sus cambios. Los filtros se aplican durante la búsqueda y no sobre los $k$ vecinos ya recuperados, lo que evita listas incompletas con filtros selectivos \citep{Gollapudi2023, Patel2024}. La búsqueda híbrida aporta la coincidencia exacta de términos poco frecuentes, como nombres de tecnologías; se eligió RRF porque no requiere calibrar puntajes ni datos de entrenamiento, con el valor habitual $k = 60$ sin ajustar, aunque \citet{Bruch2023} encontraron que RRF es sensible a ese parámetro y que una combinación convexa de los puntajes la supera. El reordenamiento se limita a 20 candidatas porque su costo crece con su número. multilingual-e5-base conserva la dimensión de e5-base-v2, de modo que el corpus se volvió a codificar en el lugar, y los vectores anteriores se conservaron para medir el efecto del cambio.

\begin{figure}[htbp!]
    \centering
    \caption{Pipeline de recuperación antes y después de la Iteración~2. En azul, lo que cambia}
    \label{fig:art-iter2}
    \begin{tikzpicture}[
        box/.style={draw, rounded corners, align=center, text width=4.4cm, minimum height=0.85cm, font=\scriptsize},
        half/.style={box, text width=3.0cm},
        new/.style={fill=AzulInstitucional!12},
        arr/.style={->, >=stealth}
    ]
        \node[font=\small\bfseries] at (2.4,0.9) {Iteración~1};
        \node[box] (a0) at (2.4,0) {Consulta};
        \node[box] (a1) at (2.4,-1.72) {Clasificación: restricciones ignoradas, o filtro SQL por subcadena del título};
        \node[box] (a2) at (2.4,-3.63) {Búsqueda densa (e5-base-v2, solo inglés): 15 fragmentos con umbral de similitud};
        \node[box] (a3) at (2.4,-5.54) {5 ofertas, ordenadas por la similitud del mejor fragmento};
        \node[box] (a4) at (2.4,-7.26) {Siempre se devuelven ofertas};
        \node[box] (a5) at (2.4,-8.98) {Generación en el idioma del contexto};
        \foreach \i/\j in {a0/a1,a1/a2,a2/a3,a3/a4,a4/a5} \draw[arr] (\i) -- (\j);
        \draw[gray!60] (6.1,-9.6) -- (6.1,1.1);
        \node[font=\small\bfseries] at (10.2,0.9) {Iteración~2};
        \node[box] (b0) at (10.2,0) {Consulta};
        \node[box, new] (b1) at (10.2,-1.72) {Clasificación: extrae salario, modalidad y tipo de contrato (inglés y español)};
        \node[half, new] (b2) at (8.45,-3.63) {Búsqueda densa (multilingual-e5-base): 50 fragmentos, con filtro};
        \node[half, new] (b3) at (11.95,-3.63) {BM25: 50 ofertas, con el mismo filtro};
        \node[box, new] (b4) at (10.2,-5.54) {Fusión RRF y reordenamiento de 20 candidatas $\rightarrow$ 5 ofertas};
        \node[box, new] (b5) at (10.2,-7.26) {Compuerta de sin resultados sobre el mejor puntaje};
        \node[box, new] (b6) at (10.2,-8.98) {Generación en el idioma de la pregunta};
        \draw[arr] (b0) -- (b1);
        \draw[arr] (b1) -- (b2); \draw[arr] (b1) -- (b3);
        \draw[arr] (b2) -- (b4); \draw[arr] (b3) -- (b4);
        \draw[arr] (b4) -- (b5); \draw[arr] (b5) -- (b6);
    \end{tikzpicture}
\end{figure}

\begin{table}[htbp!]
    \centering
    \caption{Cambios de la Iteración~2}
    \label{tab:art-cambios}
    \small
    \begin{tabular}{|p{3cm}|p{5.4cm}|p{5.4cm}|}
        \hline
        \textbf{Aspecto} & \textbf{Antes} & \textbf{Después} \\
        \hline
        Restricciones de la consulta & Ignoradas si la consulta tenía señales de recuperación; si no, filtro SQL por subcadena del título, con 30 ofertas ordenadas por fecha & Extraídas por el clasificador (inglés y español) y aplicadas como filtro dentro de la búsqueda densa y de BM25 \\
        \hline
        Primera etapa & Búsqueda densa de 15 fragmentos con umbral de similitud 0,5 & Búsqueda densa de 50 fragmentos y BM25 de 50 ofertas, fusionadas con RRF \\
        \hline
        Orden final & Similitud del mejor fragmento & Reordenamiento con Qwen3-Reranker-0.6B sobre 20 candidatas \\
        \hline
        Consultas sin ofertas relevantes & Siempre se devolvían cinco ofertas & Compuerta sobre el mejor puntaje del reordenamiento (0,134) \\
        \hline
        Modelo de \textit{embeddings} & e5-base-v2 (inglés) & multilingual-e5-base \\
        \hline
        Idioma de la respuesta & El del contexto & El de la pregunta \\
        \hline
        Salarios fuera de rango & Se almacenaban tal como llegaban & Se registran como no informados \\
        \hline
    \end{tabular}
\end{table}

\label{sec:art-rerank}

\begin{table}[!htb]
    \centering
    \caption{Selección del modelo de reordenamiento en el conjunto de desarrollo (28 consultas, 30 candidatas), con los juicios del juez local}
    \label{tab:art-rerank-dev}
    \small
    \begin{tabular}{|l|c|c|}
        \hline
        \textbf{Variante} & \textbf{nDCG@10} & \textbf{P@5} \\
        \hline
        Sin reordenamiento (densa + BM25 + RRF) & 0{,}742 & 0{,}757 \\
        \hline
        bge-reranker-v2-m3 & 0{,}804 & 0{,}814 \\
        \hline
        bge-reranker-v2-m3, fusionado con RRF & 0{,}798 & 0{,}800 \\
        \hline
        Qwen3-Reranker-0.6B & 0{,}820 & 0{,}864 \\
        \hline
        Qwen3-Reranker-0.6B con instrucción & 0{,}839 & 0{,}879 \\
        \hline
        Qwen3-Reranker-0.6B con instrucción, fusionado con RRF & 0{,}812 & 0{,}879 \\
        \hline
    \end{tabular}
\end{table}

El modelo de reordenamiento, el número de candidatas y el umbral de la compuerta se eligieron con el conjunto de desarrollo (Sección~\ref{sec:met-consultas}) y una regla de decisión fijada antes de ver los resultados, para no ajustar el sistema a las consultas de prueba; sus resultados orientan el diseño y no se reportan como resultados del sistema. Se compararon bge-reranker-v2-m3, basado en BGE-M3 \citep{Chen2024M3}, y Qwen3-Reranker-0.6B, en cinco variantes: con y sin instrucción de tarea, y con el orden dado solo por el reordenamiento o fusionado por RRF con el de la primera etapa. La regla tenía tres etapas. Primero, con 30 candidatas, se elige la variante de mayor nDCG@10, salvo que otra quede a 0,02 o menos, en cuyo caso se prefiere la más simple. Después, se usan 20 candidatas en lugar de 30 si el nDCG@10 baja 0,01 o menos. Por último, el umbral de la compuerta es el mayor valor que deja por debajo, como máximo, a una de las 28 consultas de desarrollo, evaluado sobre el mejor puntaje de cada consulta, y se adopta solo si deja por debajo al menos a 8 de las 10 consultas fuera de dominio.


Todas las variantes mejoraron entre 0,06 y 0,10 el nDCG@10 de la recuperación sin reordenamiento (Tabla~\ref{tab:art-rerank-dev}). Qwen3-Reranker-0.6B sin instrucción quedó a 0,019 de la mejor variante, y la regla lo eligió por ser más simple. Pasar de 30 a 20 candidatas cambió el nDCG@10 de 0,832 a 0,830 y redujo la mediana de la latencia de recuperación de 1.184 a 788~ms, así que se adoptaron 20. El 0,832 de la variante con 30 candidatas no coincide con el 0,820 de la Tabla~\ref{tab:art-rerank-dev} porque cada etapa de la regla se juzgó con su propio agrupamiento (Sección~\ref{sec:met-juicios}). Las listas de la variante son idénticas en las dos etapas y sus 540 pares comunes tienen las mismas etiquetas, pero la etapa~1 agrupó las listas de seis variantes, con 15,0 ofertas relevantes por consulta, y la etapa~2 las de tres, con 13,7. Con menos ofertas relevantes, el ordenamiento ideal con que se normaliza el nDCG es menor, y la misma lista obtiene un valor más alto. Las dos cantidades de candidatas se compararon dentro de la etapa~2, con el mismo agrupamiento, de modo que la decisión no depende de esa diferencia. El mejor puntaje más bajo entre las consultas de desarrollo fue 0,132 y el siguiente 0,134; las diez consultas fuera de dominio quedaron por debajo de 0,09, de modo que el umbral de 0,134 cumplió la condición y se adoptó.

Los juicios de relevancia de esta selección los dio el juez local (Llama 3.1 8B), el único disponible entonces, cuyo acuerdo con las etiquetas humanas resultó después moderado ($\kappa$ binario de \cKappaLocal, Sección~\ref{sec:res-juez}). La regla se aplicó, por tanto, con el juez de desarrollo. El aporte del reordenamiento se confirmó después en el conjunto de prueba con el juez final (\cDifReord{} de nDCG@10, Sección~\ref{sec:res-pe1}).
